\section{A fully opened literal mass bound}
\label{sec:mass-bound}

\subsection{Why the moving ultra sum must be opened first}

On the left-ultra polarization the signed factor
\[
 C_{m_\alpha}(j)
 =
 \sum_{\substack{u\mid N_\alpha(j)\\T_X<u\le U_{0,X}}}
 a_R(u)
\]
contains both conductor-one children
\[
 q_X\mid\sigma=\frac{N_\alpha(j)}u
\]
and nonconstant children.  Cancellation between those two classes
can make \(C_{m_\alpha}(j)\) smaller than the conductor-one
subsum.  Therefore \(|C_{m_\alpha}(j)|\) is not an upper bound for
the absolute constant-sector child mass.  The row-gcd projector has
the same issue: its children may cancel even when their absolute
constant-sector mass is nonzero.

We consequently take total variation only after every relevant child
has been opened.  This order also includes projector children whose
unexpanded parent is killed by the actual mask; such children cancel
in the full projector identity but cannot be silently deleted from a
sectorwise absolute majorant.

\subsection{Columnwise projector variation}

Fix an opposite row
\(\gamma=(\ell_\gamma,e_\gamma)\).  Let
\(\lambda_{G_X}(v)\) be the exact row-gcd projector coefficient.
If the local coprimality indicator is also expanded, its coefficient
mass is at most \(2^{\omega(v)}\).  Put
\begin{equation}
 \mathfrak P_\gamma
 =
 \sum_{v\mid e_\gamma}
 |\lambda_{G_X}(v)|\,2^{\omega(v)}.
\label{eq:column-projector-variation}
\end{equation}
The literal TPC-33 divisor bounds give, uniformly in the fixed
column,
\begin{equation}
 \mathfrak P_\gamma\le X^{o(1)}.
\label{eq:column-projector-soft}
\end{equation}
For example, on squarefree \(e_\gamma\), the elementary bounds
\(|\lambda_{G_X}(v)|\le\tau(v)\) and
\(\tau(v)=2^{\omega(v)}\) give
\(\mathfrak P_\gamma\le5^{\omega(e_\gamma)}=X^{o(1)}\).

The point is the order of quantifiers: \eqref{eq:column-projector-soft}
is used inside the sum over \(\gamma\).  It is never promoted to a
low-rank decomposition of the full row matrix.

\subsection{The fully opened left majorant}

Let
\[
 \mathcal U_q(\alpha,j)
 =
 \left\{
 u:
 \begin{array}{c}
 u\mid N_\alpha(j),\quad T_X<u\le U_{0,X},\\
 q_X\mid N_\alpha(j)/u
 \end{array}
 \right\}.
\]
After opening the ultra divisor, row-gcd projector, local
coprimality children, exact-content children, finite residue labels,
and interval components, the total variation of the left
conductor-one sector is bounded by
\begin{align}
 \mathcal V_L
 \le{}&
 X^{o(1)}
 \sum_\gamma|\gamma_\gamma^{(2)}|
 \mathfrak P_\gamma
 \sum_\alpha|\gamma_\alpha^{(1)}|
 \sum_{\substack{j\in\cJ_X\\
                 q_X\mid N_\alpha(j)}}
 \notag\\
 &\qquad\times
 \left(
 \sum_{u\in\mathcal U_q(\alpha,j)}
 |a_R(u)|
 \right)
 |\mathsf H_{m_\gamma}(j)|.
\label{eq:fully-opened-left-majorant}
\end{align}
Here the leading soft factor includes:
\begin{itemize}[leftmargin=2em]
\item the exact-content total variation
  \[
  \sum_{\substack{c\mid G\\c\le C_X}}
  \sum_{\kappa\mid G/c}|\mu(\kappa)|
  \le\tau(G)^2=X^{o(1)};
  \]
\item at most two interval components and the finite periodic/residue
  separation; and
\item the total-variation norm of the absolutely convergent
  Mellin/BV representation.
\end{itemize}
The finite projective mass and the uniform Mellin/BV
total-variation control are precisely the separation statements
proved in TPC-25 and TPC-33, uniformly in
\((\alpha,\gamma,j,u,v,c,\kappa)\)
\citep{WangTPC25,WangTPC33,WangTPC93}.
Tonelli's theorem applies to this nonnegative majorant.  Thus
continuous separation variables remain inside their total-variation
integral and do not become fictitious finite native keys.

\begin{proposition}[Fully opened polarization bound]
\label{prop:polarization-bound}
Each exact polarization has conductor-one child variation
\[
 \mathcal V_L,\mathcal V_R
 \ll_{\mathscr D,h_0}X^{o(1)}Q_X^2.
\]
\end{proposition}

\begin{proof}
For each moving row \(\alpha\), the \(j\)-sum in
\eqref{eq:fully-opened-left-majorant} contains at most one point by
\cref{thm:orbit-sparsity}.  At that point,
\[
 \sum_{u\in\mathcal U_q(\alpha,j)}|a_R(u)|
 \le
 \sum_{\substack{u\mid N_\alpha(j)\\T_X<u\le U_{0,X}}}
 |\mu(u)|\log u
 \le X^{o(1)}
\]
by \eqref{eq:ultra-envelope}, while
\(|\mathsf H_{m_\gamma}(j)|\le X^{o(1)}\) by
\eqref{eq:prefix-envelope}.  Hence
\begin{align*}
 \mathcal V_L
 &\le
 X^{o(1)}
 \sum_\gamma|\gamma_\gamma^{(2)}|\mathfrak P_\gamma
 \sum_\alpha|\gamma_\alpha^{(1)}|\\
 &\le
 X^{o(1)}
 \left(\sum_\gamma|\gamma_\gamma^{(2)}|\right)
 \left(\sum_\alpha|\gamma_\alpha^{(1)}|\right)
 \ll X^{o(1)}Q_X^2
\end{align*}
by \eqref{eq:column-projector-soft} and \eqref{eq:row-l1}.
The right-ultra term is identical after exchanging the rows.
\end{proof}

\subsection{Restoring the reconstruction frequencies}

For fixed \(K\ge2\), the translated low-window reconstruction
multiplier obeys
\begin{equation}
 \sum_{0<d_{q_X}(r,0)\le R_X}
 |\mathfrak m_{K,X}(r)|
 \ll_K1.
\label{eq:filter-l1}
\end{equation}
The conductor-one predicate is independent of the nonzero low
frequency after the literal divisibility \(c_\xi\mid\Omega_\xi\) has
been imposed.  The orientation and fixed phase have modulus one.
Thus restoring all frequencies costs only
\eqref{eq:filter-l1}, not their cardinality.

\begin{theorem}[Endpoint-scale constant-sector bound]
\label{thm:literal-constant-bound}
Let \(\mathcal L^{\rm const}_{K,R,X}\) be the sum of all fully opened
literal TPC-93 low-window children that
\begin{enumerate}[label=\textup{(\roman*)}]
\item occur in either exact polarization, including children of
  projector parents that cancel only after reassembly;
\item retain every native mask, exact-content label, interval
  component, prefix, and original coefficient; and
\item have algebraic conductor \(1\), equivalently
  \(q_X\mid\Omega_\xi\).
\end{enumerate}
Then
\begin{equation}
 \boxed{
 |\mathcal L^{\rm const}_{K,R,X}|
 \ll_{K,h_0,\mathscr D}
 X^{o(1)}Q_X^2.}
\label{eq:main-constant-bound}
\end{equation}
\end{theorem}

\begin{proof}
The absolute value of the sector sum is at most
\(\mathcal V_L+\mathcal V_R\) times the frequency multiplier
\(\ell^1\) mass.  Apply
\cref{prop:polarization-bound} and \eqref{eq:filter-l1}.
\end{proof}

\begin{remark}[No quotient estimate]
The proof works directly with opened physical children, fixed
opposite columns, and their native total variation.  Equal
determinants, affine forms, phase multipliers, or targets are never
merged.  Neither signed prefix cancellation nor projector
cancellation is used as an upper bound for a selected sector.
\end{remark}
